% In this file you should put the actual content of the blueprint.
% It will be used both by the web and the print version.
% It should *not* include the \begin{document}

% This makes the macros `\inputleanmodule` and `\inputleannode` available.
% The statements and proofs of all numbered results live in the Lean sources, as
% `@[blueprint (statement := ...) (proof := ...)]` annotations, and are extracted by
% `lake build :blueprint`.  This file only supplies the narrative and the ordering.
\input{../../.lake/build/blueprint/library/BV}

\chapter{Introduction}

Let $\Lambda$ be the von Mangoldt function and, for a modulus $q \ge 1$ and a residue class
$a \bmod q$, let
$$\psi(x; q, a) := \sum_{\substack{n \le x \\ n \equiv a \pmod q}} \Lambda(n)$$
be Chebyshev's function in the arithmetic progression $a \bmod q$. The prime number theorem in
arithmetic progressions predicts $\psi(x; q, a) \approx x/\varphi(q)$ for $(a, q) = 1$, but the
best unconditional estimates for an \emph{individual} modulus $q$ (the Siegel--Walfisz theorem)
only cover $q \le (\log x)^C$, and even the generalised Riemann hypothesis would only give
$q \le x^{1/2 - \varepsilon}$. The Bombieri--Vinogradov theorem shows that the latter range is
available \emph{on average} over the modulus: for every $A \ge 0$,
$$\sum_{q \le Q} \max_{y \le x} \max_{a \in (\Z/q\Z)^*} \left| \psi(y; q, a) - \frac{y}{\varphi(q)} \right| \ll_A \frac{x}{(\log x)^A}
\qquad \text{uniformly for } x \ge 2 \text{ and } 1 \le Q \le \frac{\sqrt x}{(\log x)^{A+3}} .$$
This is often described as ``the generalised Riemann hypothesis on average'', and it is the
main analytic input in sieve-theoretic applications such as the work of Goldston--Pintz--Y\i ld\i r\i m
and Maynard on small gaps between primes.

This blueprint describes a formalisation of the Bombieri--Vinogradov theorem in Lean~4 and
Mathlib, following Chapter~26 of Koukoulopoulos' \emph{The Distribution of Prime Numbers}. Two
analytic results are taken as external inputs (axioms in Lean): the Siegel--Walfisz theorem
(\Cref{siegel_walfisz}) and the multiplicative large sieve inequality (\Cref{large_sieve}). Everything
else is proved from Mathlib and the PrimeNumberTheoremAnd library
(from which we borrow the smoothed cutoff function used in the Perron-type argument of
\Cref{ch:typeII}).

The text has been written so as to match the Lean development exactly: every numbered statement
carries the hypotheses that its Lean counterpart carries, and each proof follows the steps of the
formal proof. Implied constants are not tracked: in Lean every $\ll$ below is witnessed by some real
number, obtained where convenient by choice or compactness arguments, and nothing is claimed about its
size (in particular the implied constants in the Siegel--Walfisz axiom are not effective). Several statements therefore differ from the textbook presentation; the most important
deviations are recorded in \Cref{ch:overview}.

\begin{theorem*}[Main result, informal form]
Assume the Siegel--Walfisz theorem and the large sieve inequality. Then for every $A \in \N$, all real
$x \ge 2$ and $1 \le Q \le \sqrt x/(\log x)^{A+3}$,
$$\sum_{q \le Q}\ \sup_{1 \le y \le x}\ \max_{a \in (\Z/q\Z)^*} \left| \psi(y; q, a) - \frac{y}{\varphi(q)} \right| \ll_A \frac{x}{(\log x)^A} .$$
\end{theorem*}
The precise statement is \Cref{Bombieri-Vinogradov} (with a real-valued corollary
\Cref{bombieri_vinogradov}).

\chapter{Overview of the proof}\label{ch:overview}

\section{Discrepancies and the reduction to \texorpdfstring{$\Lambda$}{Lambda}}
For an arithmetic function $f$ we measure its distribution in progressions by the
\emph{discrepancy}
$$\Delta_f(y; q, a) := \sum_{\substack{n \le y \\ n \equiv a \pmod q}} f(n) - \frac{1}{\varphi(q)} \sum_{\substack{n \le y \\ (n,q) = 1}} f(n)$$
(\Cref{Delta}). For $f = \Lambda$ the difference between $\psi(y; q, a) - y/\varphi(q)$ and
$\Delta_\Lambda(y; q, a)$ is $\varphi(q)^{-1}\big(\sum_{n \le y, (n,q)=1} \Lambda(n) - y\big)$, which is
controlled by the prime number theorem with logarithmic savings (a special case of Siegel--Walfisz,
\Cref{PNT}) together with the elementary bound $\sum_{n \le y, (n,q) > 1} \Lambda(n) \le \log q \log y/\log 2$
(\Cref{sum_vonMangoldt_not_coprime_le}). Summing $1/\varphi(q)$ over $q \le Q$ costs a factor $\log x$
(\Cref{summatory_totient_inv_le}). The main theorem is thus reduced to the statement
\Cref{BV_Delta_Lambda_enorm} about $\Delta_\Lambda$, in which the maximum over $y$ is restricted to
$\sqrt x \le y \le x$; the range $y < \sqrt x$ is handled by the trivial bound
$|\psi(y;q,a) - y/\varphi(q)| \ll y$ (\Cref{Chebyshev.abs_psiMod_sub_div_le_const}).

\section{Vaughan's identity}
Fix parameters $U, V$ with $UV \le \sqrt x$ and $U, V \ge e^{\sqrt{\log x}}$ (\Cref{ProofData}). Vaughan's
identity (\Cref{Lambda_decomp}) decomposes
$$\Lambda = \Lambda^\sharp + \Lambda^\flat + \Lambda_{\le U}, \qquad
\Lambda^\sharp = \mu_{\le V} * \log - \Lambda_{\le U} * \mu_{\le V} * 1, \qquad
\Lambda^\flat = (\Lambda_{> U} * 1) * \mu_{> V} .$$
The three pieces are treated separately (\Cref{ch:small,ch:typeI,ch:typeII}) and
recombined by the triangle inequality in \Cref{ch:assembly}. The piece $\Lambda_{\le U}$ is
supported on $[1, U]$ and is bounded trivially. The piece $\Lambda^\sharp$ is a \emph{Type~I} sum: a
convolution of a coefficient function of small support with the smooth functions $\log$ and $1$,
whose discrepancies are bounded pointwise by Abel summation and $|\Delta_1| \le 1$. The piece
$\Lambda^\flat$ is a \emph{Type~II} (bilinear) sum, which is bounded on average over $q$ by expanding
$\Delta_{\Lambda^\flat}$ in Dirichlet characters, grouping characters by conductor, treating the
small conductors with Siegel--Walfisz, and treating the large conductors with the large sieve.

\section{The two external inputs}
\emph{Siegel--Walfisz} (\Cref{siegel_walfisz}) is used in three places: with $q = 1$ as the prime number
theorem (\Cref{PNT}), in the reduction from $\psi$ to $\Delta_\Lambda$, and in the treatment of the
small conductors $d \le (\log x)^C$ in the Type~II sum (\Cref{Delta_onCoprime_Lambda_bound}). The
\emph{large sieve} (\Cref{large_sieve}) enters only through the maximal large sieve inequality for
convolutions (\Cref{LargeSieve_convolution}, Theorem~26.6 of Koukoulopoulos), which we prove via a
smoothed Perron formula rather than the truncated Perron formula of the textbook (\Cref{ch:typeII}).

\section{Where the formalisation deviates from the textbook}\label{sec:deviations}
The following points are worth keeping in mind when comparing with Koukoulopoulos.
\begin{itemize}
\item \emph{A divisor factor in the Type~I bound.} The pointwise bound for the restricted function
$\Lambda^\sharp_r$ carries a factor $\tau(r)$ (\Cref{Delta_LambdaSharp_bound}): restriction to integers
coprime to $r$ commutes with convolution (\Cref{ArithmeticFunction.on_mul_of_saturated}), so the smooth
factor $\log$ becomes $\log_r$, which is neither smooth nor monotone, and we Möbius-expand it as a sum of
$\tau(r)$ dilated copies of $\log$ and $1$ (\Cref{zeta_on_coprime_apply}). The factor is harmless: in the
Type~I average it is summed over the modulus using $\sum_{q \le Q} \tau(q) \le 3 Q \log Q$
(\Cref{sum_divisors_card_le}), and in the small-conductor analysis it is absorbed by $\tau(q) \le 2\sqrt q \le 2x^{1/4}$,
which is why the small-conductor lemmas assume $q \le \sqrt x$.
\item \emph{The large sieve for convolutions is proved by smoothing.} Instead of the truncated Perron
formula, we use the smoothed cutoff $\widetilde 1_\varepsilon$ of the PrimeNumberTheoremAnd library
(\Cref{Flat.Bump}). At half-integers $y = K + \tfrac12$ the smoothed and sharp cutoffs agree exactly
when $\varepsilon = 1/(6 \log 2 \cdot x)$ (\Cref{Flat.T_eq_sharp}), so no error term arises, and Mellin
inversion on the line $\Re s = 1/\log(x+1)$ (\Cref{Flat.T_mellin_representation}) reduces the problem to
the large sieve for Dirichlet polynomials (\Cref{Flat.largeSieve_char_bound}) and to a kernel integral
of size $O(\log x)$ (\Cref{Flat.mellin_J_bound}).
\item \emph{Uniformity in $U, V$.} All Type~I/II estimates are proved for an arbitrary parameter datum
(\Cref{ProofData}); the choice $U = V = e^{\sqrt{\log x}}$ is only made in the final assembly. The
constraints force $\log x \ge 16$ (\Cref{log_x_large}), so the main theorem splits into the range
$x \le e^{32}$, where the trivial bound suffices (\Cref{BV_compact_enorm}), and the range $x > e^{32}$.
\item \emph{Maxima are extended-real suprema.} The maxima over $y$ and $a$ are formalised as suprema in
$[0, \infty]$ (\Cref{BV.maxy}, \Cref{BV.maxya}), so that no boundedness hypotheses need to be carried
around; finiteness is part of each estimate.
\end{itemize}

\chapter{Definitions and preliminaries}\label{ch:prelim}

\section{Notation}
$\N = \{0, 1, 2, \dots\}$ denotes the natural numbers (as in Lean), although all sums over $n \le x$
start at $n = 1$. For $a, b \in \N$ we write $(a, b) = \gcd(a, b)$ and $\tau(n)$ for the number of
divisors of $n$. Dirichlet convolution is written $f * g$, the constant function $1$ on $\N_{\ge 1}$
(Mathlib's $\zeta$) is written $1$, $\mu$ is the Möbius function, $\varphi$ Euler's totient and
$\Lambda$ the von Mangoldt function; $\log$ also denotes the arithmetic function $n \mapsto \log n$
and $\log^v$ its pointwise $v$-th power (with $\log^0 = 1$). For a Dirichlet character $\chi$
modulo $q$ we write $\chi_0$ for the principal character, $\sumstar_{\chi \bmod q}$ for a sum over
primitive characters, and $\xi \uparrow^q$ for the character modulo $q$ induced by a character
$\xi$ modulo a divisor of $q$. Complex conjugation is $\overline{z}$.

For real-valued quantities $f, g$ the statement $f \ll g$ (equivalently $f = O(g)$) means $|f| \le C g$ for
some constant $C > 0$; a subscript, as in $f \ll_A g$, records the parameters the constant may depend on.
Unless stated otherwise the constant may also depend on the implied constants of the two external inputs
(\Cref{siegel_walfisz}, \Cref{large_sieve}) and on the fixed bump function of \Cref{Flat.Bump}, but never
on $x, y, q, a, Q$ or on the parameter datum. Constants are not tracked in the formalisation.

\section{Summatory functions and Chebyshev's function}
\inputleannode{summatory}
\inputleannode{Nat.modEqs}
\inputleannode{Chebyshev.psiMod}

\section{The discrepancy \texorpdfstring{$\Delta_f$}{Delta\_f}}
\inputleannode{onCoprime}
\inputleannode{ArithmeticFunction.on}
\inputleannode{dilate}
\inputleannode{Delta}
\inputleannode{Delta_linear}
\inputleannode{abs_Delta_le_two_summatory_abs}
\inputleannode{Delta_onCoprime_self}
\inputleannode{Delta_eq_sum_char}
\inputleannode{Delta_Lambda_eq}

\section{The parameters}
\inputleannode{ProofData}
\inputleannode{log_x_large}
\inputleannode{UV_range}

\section{Maximal quantities}
\inputleannode{BV.maxy}
\inputleannode{BV.maxya}
\inputleannode{maxya_basic}
\inputleannode{maxya_Delta_enorm_le_of_abs_le}
\inputleannode{maxya_Delta_add_le}

\chapter{External analytic inputs}\label{ch:axioms}

Two analytic results are quoted from the literature and taken as axioms in Lean. In a proof,
``external input'' means that the assertion is quoted and no proof is given here. Both are
stated with an unspecified constant, which is itself part of the axiom.

\section{The Siegel--Walfisz theorem}
\inputleannode{siegel_walfisz}
\inputleannode{PNT}
\inputleannode{sum_vonMangoldt_not_coprime_le}
\inputleannode{coprime_vonMangoldt_error}
\inputleannode{pnt_ratio_bound}

\section{The large sieve inequality}
\inputleannode{large_sieve}

\chapter{Auxiliary arithmetic estimates}\label{ch:aux}

This chapter collects the elementary estimates used in the main argument: the trivial bound for
the modular Chebyshev error, the average order of the divisor function, the convergent sums
involving $1/\varphi$, and the comparisons between powers of $\log x$ and powers of $x$.

\inputleannode{Chebyshev.abs_psiMod_sub_div_le_const}
\inputleannode{sum_divisors_card_le}
\inputleannode{card_divisors_le_two_mul_sqrt}
\inputleannode{d_le_two_mul_totient_sq}
\inputleannode{invTotientAF_eq_mul}
\inputleannode{summatory_totient_inv_le}
\inputleannode{sum_rphiInv_le}
\inputleannode{log_pow_le_const_mul_rpow}
\inputleannode{log_power_absorption}
\inputleannode{log_pow_le_const_mul_sqrt}

\chapter{Vaughan's identity}\label{ch:vaughan}

Throughout the rest of the blueprint a parameter datum $(x, U, V)$ (\Cref{ProofData}) is fixed.

\inputleannode{LambdaLEU}
\inputleannode{moebiusLEV}
\inputleannode{LambdaSharp}
\inputleannode{LambdaFlat}
\inputleannode{Lambda_decomp}

\chapter{The small part \texorpdfstring{$\Lambda_{\le U}$}{Lambda le U}}\label{ch:small}

The contribution of $\Lambda_{\le U}$ is bounded trivially: it is supported on $[1, U]$ with
$U \le \sqrt x$, and there are at most $Q \le \sqrt x/(\log x)^{A+3}$ moduli.

\inputleannode{sum_LambdaLEU_le}
\inputleannode{Delta_LambdaLEU_bound}
\inputleannode{BV_LambdaLE_enorm}

\chapter{Type I sums: \texorpdfstring{$\Lambda^\sharp$}{Lambda sharp}}\label{ch:typeI}

\section{Discrepancy of smooth weights}
The basic Type~I mechanism: for $f$ of small support, $\Delta_{f * g}$ is a short sum of
discrepancies of $g$ at the points $y/k$ (\Cref{Delta_convolution_eq}); when $g = \log^v$ these are bounded
pointwise via Abel summation against the trivial bound $|\Delta_1| \le 1$.
\inputleannode{Delta_convolution_eq}
\inputleannode{Delta_one_bound}
\inputleannode{Delta_abel_summation}
\inputleannode{Delta_monotone_bound}
\inputleannode{Delta_flog_bound}

\section{Restriction to coprime integers and Möbius expansion}
We need bounds not for $\Lambda^\sharp$ itself but for its restriction $\Lambda^\sharp_r$ to integers coprime
to $r$ (this is what the conductor analysis of \Cref{ch:typeII} produces). Restriction is a
convolution homomorphism, so the smooth factor $\log$ becomes the restricted function $\log_r$, to
which \Cref{Delta_flog_bound} does not apply directly. We rewrite $\log_r$ (and $1_r$) as a Möbius sum of
dilations of the unrestricted functions; this costs a factor $\tau(r)$.
\inputleannode{ArithmeticFunction.on_mul_of_saturated}
\inputleannode{mul_dilate}
\inputleannode{summatory_abs_dilate_le}
\inputleannode{summatory_abs_mul_le}
\inputleannode{moebius_coprime_indicator}
\inputleannode{zeta_on_coprime_apply}
\inputleannode{mul_on_coprime_expansion}
\inputleannode{Delta_dilate_flog_bound}

\section{The pointwise bound for \texorpdfstring{$\Lambda^\sharp_r$}{Lambda sharp r}}
\inputleannode{l1_norms_coefficients}
\inputleannode{Delta_term1_bound}
\inputleannode{Delta_term2_bound}
\inputleannode{Delta_LambdaSharp_bound}

\section{Summing over the modulus}
\inputleannode{BV_LambdaSharp_enorm}

\chapter{Type II sums: \texorpdfstring{$\Lambda^\flat$}{Lambda flat}}\label{ch:typeII}

\section{Definitions}
\inputleannode{S}
\inputleannode{T}

\section{Characters and conductors}
By orthogonality, $\varphi(q)\Delta_{\Lambda^\flat}(y;q,a)$ is a sum over non-principal characters
$\chi \bmod q$. Grouping the characters by conductor $d \mid q$ produces, for each $d$, a sum over the
primitive characters modulo $d$ of the twisted sums $S_{q/d}(y, \xi)$. The conductors
$d \le (\log x)^C$ are treated in \Cref{sec:small-cond} by Siegel--Walfisz, after converting the
primitive-character sums back into discrepancies by Möbius inversion; the conductors
$d > (\log x)^C$ are collected into the quantities $T_r$ and treated by the large sieve.
\inputleannode{character_sum_by_conductor}
\inputleannode{totient_mul_Delta_eq}
\inputleannode{character_sum_Mobius}
\inputleannode{Delta_LambdaFlat_decomp}

\section{Small conductors}\label{sec:small-cond}
\inputleannode{Delta_onCoprime_Lambda_bound}
\inputleannode{Delta_onCoprime_LambdaFlat_pointwise}
\inputleannode{Delta_LambdaFlat_small_conductor}

\section{Reduction to the quantities \texorpdfstring{$T_r$}{T\_r}}
\inputleannode{maxya_Delta_LambdaFlat_enorm_le}
\inputleannode{BV_LambdaFlat_via_T}

\section{A smoothed Perron formula}\label{sec:perron}
The remaining input is a maximal version of the large sieve for Dirichlet convolutions
(\Cref{LargeSieve_convolution}). The textbook proof uses the truncated Perron formula; we use instead
a smoothed cutoff $\widetilde 1_\varepsilon$, whose Mellin transform is $s^{-1}\mathcal M[\nu](\varepsilon s)$ for a fixed
bump function $\nu$. The point is that for $y = K + \tfrac12$ a half-integer and
$\varepsilon = 1/(6\log 2 \cdot x)$, the smoothed cutoff coincides with the sharp cutoff at every integer
(\Cref{Flat.T_eq_sharp}), so the smoothed bilinear sums $T_\varepsilon(y, \chi)$ are \emph{equal} to the
partial sums we want to bound, and no error term needs to be estimated. In this section $x \ge 1$ and
$Q \ge 1$ are arbitrary reals, $(f, g)$ is a supported pair and $\nu$ a bump function; no parameter
datum is involved.
\subsection*{Inputs from the PrimeNumberTheoremAnd library}
The smoothed cutoff and its basic properties are taken from the PrimeNumberTheoremAnd library
(module \texttt{MellinCalculus}). We record the statements we use; the proofs are in that library and
are not reproduced here. Throughout, $\nu : \R \to \R$ is a $C^1$ function supported in $[1/2, 2]$.

\begin{theorem}[Existence of a bump function]\label{SmoothExistence}\lean{SmoothExistence}\leanok
There exists a $C^\infty$ function $\nu : \R \to \R$, supported in $[1/2, 2]$, with $\nu \ge 0$ everywhere and
$\int_0^\infty \nu(t)\, \frac{\mathrm{d}t}{t} = 1$.
\end{theorem}
\begin{proof}\leanok Proved in PrimeNumberTheoremAnd. \end{proof}

\begin{definition}[The smoothed cutoff $\widetilde 1_\varepsilon$]\label{Smooth1}\lean{Smooth1}\leanok
For $\varepsilon > 0$, the smoothed cutoff $\widetilde{1}_\varepsilon : \R \to \R$ is the multiplicative
(Mellin) convolution of the indicator of $(0, 1]$ with the spike $t \mapsto \nu(t^{1/\varepsilon})/\varepsilon$:
$$\widetilde{1}_\varepsilon(u) := \int_0^\infty 1_{(0,1]}(t)\, \frac{1}{\varepsilon}\, \nu\big((u/t)^{1/\varepsilon}\big)\, \frac{\mathrm{d}t}{t} .$$
\end{definition}

\begin{lemma}[The smoothed cutoff equals $1$ below $1 - (\log 2)\varepsilon$]\label{Smooth1Properties_below}\lean{Smooth1Properties_below}\leanok
\uses{Smooth1}
Assume $\int_0^\infty \nu(t)\,\mathrm{d}t/t = 1$. For $\varepsilon > 0$ and $0 < u \le 1 - (\log 2)\varepsilon$ we have
$\widetilde 1_\varepsilon(u) = 1$.
\end{lemma}
\begin{proof}\leanok Proved in PrimeNumberTheoremAnd. \end{proof}

\begin{lemma}[The smoothed cutoff vanishes above $1 + 2(\log 2)\varepsilon$]\label{Smooth1Properties_above}\lean{Smooth1Properties_above}\leanok
\uses{Smooth1}
For $0 < \varepsilon < 1$ and $u \ge 1 + 2(\log 2)\varepsilon$ we have $\widetilde 1_\varepsilon(u) = 0$.
\end{lemma}
\begin{proof}\leanok Proved in PrimeNumberTheoremAnd. \end{proof}

\begin{lemma}[Non-negativity of the smoothed cutoff]\label{Smooth1Nonneg}\lean{Smooth1Nonneg}\leanok
\uses{Smooth1}
Assume $\nu \ge 0$ on $(0, \infty)$. Then $\widetilde 1_\varepsilon(u) \ge 0$ for all $\varepsilon > 0$ and $u > 0$.
\end{lemma}
\begin{proof}\leanok Proved in PrimeNumberTheoremAnd. \end{proof}

\begin{lemma}[The smoothed cutoff is at most $1$]\label{Smooth1LeOne}\lean{Smooth1LeOne}\leanok
\uses{Smooth1}
Assume $\nu \ge 0$ on $(0, \infty)$ and $\int_0^\infty \nu(t)\,\mathrm{d}t/t = 1$. Then $\widetilde 1_\varepsilon(u) \le 1$
for all $\varepsilon > 0$ and $u > 0$.
\end{lemma}
\begin{proof}\leanok Proved in PrimeNumberTheoremAnd. \end{proof}

\begin{lemma}[Continuity of the smoothed cutoff]\label{Smooth1ContinuousAt}\lean{Smooth1ContinuousAt}\leanok
\uses{Smooth1}
Assume $\nu \ge 0$ on $(0,\infty)$. For $\varepsilon > 0$, $\widetilde 1_\varepsilon$ is continuous at every $u > 0$.
\end{lemma}
\begin{proof}\leanok Proved in PrimeNumberTheoremAnd. \end{proof}

\begin{lemma}[Mellin transform of the smoothed cutoff]\label{MellinOfSmooth1a}\lean{MellinOfSmooth1a}\leanok
\uses{Smooth1}
For $\varepsilon > 0$ and $\Re s > 0$,
$$\mathcal{M}[\widetilde 1_\varepsilon](s) = \frac{1}{s}\, \mathcal{M}[\nu](\varepsilon s), \qquad \text{where } \mathcal M[f](s) := \int_0^\infty u^{s-1} f(u)\,\mathrm{d}u .$$
\end{lemma}
\begin{proof}\leanok Proved in PrimeNumberTheoremAnd. \end{proof}

\begin{lemma}[Decay of the Mellin transform of the smoothed cutoff]\label{MellinOfSmooth1b}\lean{MellinOfSmooth1b}\leanok
\uses{Smooth1, MellinOfSmooth1a}
There is a constant $C > 0$, depending only on $\nu$, such that for all $\sigma_1 > 0$, all $s$ with
$\sigma_1 \le \Re s \le 2$ and all $0 < \varepsilon < 1$,
$$\|\mathcal{M}[\widetilde 1_\varepsilon](s)\| \le \frac{C}{\varepsilon \|s\|^2} .$$
\end{lemma}
\begin{proof}\leanok Proved in PrimeNumberTheoremAnd. \end{proof}

\begin{lemma}[Integration by parts on $(0,\infty)$]\label{PartialIntegration}\lean{PartialIntegration}\leanok
Let $f, g : \R \to \C$ be differentiable on $(0, \infty)$ with $f g'$ and $f' g$ integrable on $(0,\infty)$ and
$f(u)g(u) \to 0$ as $u \to 0^+$ and as $u \to \infty$. Then $\int_0^\infty f g' = -\int_0^\infty f' g$.
\end{lemma}
\begin{proof}\leanok Proved in PrimeNumberTheoremAnd. \end{proof}

\subsection*{The smoothed bilinear sums}
\inputleannode{Flat.Bump}
\inputleannode{Flat.FG}
\inputleannode{Flat.T}
\inputleannode{mellin_bump_bounded}
\inputleannode{Smooth1_mellinInv_mellin_eq}
\inputleannode{mellin_smooth1_bounds}
\inputleannode{Flat.mellin_J_bound}
\inputleannode{Flat.sup_summatory_eq_sup_nat}
\inputleannode{Flat.FG.summatory_mul_char}
\inputleannode{Flat.T_eq_sharp}
\inputleannode{Flat.T_mellin_representation}

\section{The large sieve for convolutions}
\inputleannode{Flat.largeSieve_factor}
\inputleannode{Flat.largeSieve_char_bound}
\inputleannode{Flat.summatory_T_ll}
\inputleannode{Flat.LargeSieve_convolution_aux}
\inputleannode{LargeSieve_convolution}

\section{Dyadic decomposition and the bound for \texorpdfstring{$T_r$}{T\_r}}
\inputleannode{LambdaFlat_dyadic}
\inputleannode{BV_char_sum_bound}
\inputleannode{T_r_bound}

\section{The Type II estimate}
\inputleannode{BV_LambdaFlat_enorm}

\chapter{Assembly}\label{ch:assembly}

\inputleannode{maxya_Delta_Lambda_enorm_le}
\inputleannode{BV_Delta_Lambda_enorm}
\inputleannode{psiMaxEnorm}
\inputleannode{BV_compact_enorm}
\inputleannode{BV_large_enorm}
\inputleannode{Bombieri-Vinogradov}
\inputleannode{bombieri_vinogradov}
